% Grenzwerte und Verhalten im Unendlichen – bank/grenzwerte-und-verhalten-im-unendlichen/ (zusammenbau v0.4, Heft)
\einheitenkopf[t6]{Grenzwerte und Verhalten im Unendlichen}
\setcounter{aufgabe}{54}

% Anlauf (ohne Prüfkennung)
\begin{aufgabe}{Ganzrationale Funktionen – Anlauf}
\begin{teile}
\teil Gerade oder ungerade, plus oder minus? $f(x) = -3x^5 + x^2$ \kreuz{Grad gerade, Koeffizient positiv} \\ \kreuz{Grad gerade, Koeffizient negativ} \\ \kreuz{Grad ungerade, Koeffizient positiv} \\ \kreuz{Grad ungerade, Koeffizient negativ}
\teil $f(x) = x^2 - 6x + 1$. Verhalten für $x \to +\infty$ und $x \to -\infty$? x → +∞: f(x) → \leerfeld ; x → −∞: f(x) → \leerfeld
\end{teile}
\end{aufgabe}

\begin{aufgabe}{Ganzrationale Funktionen – Prüfungsaufgaben}
\begin{teile}
\teil $f(x) = -\frac{1}{3}x^3 + 2x^2$. Gib das Verhalten der Funktionswerte für $x \to +\infty$ und $x \to -\infty$ an. \hfill (Abitur 2022 GK) \\ x → +∞: f(x) → \leerfeld ; x → −∞: f(x) → \leerfeld
\teil $f(x) = \frac{1}{20}x^5 - x^2$ hat die Ableitung $f'(x) = \frac{1}{4}x^4 - 2x$. Gib das Verhalten von $f'$ für $x \to -\infty$ und $x \to +\infty$ an. \hfill (Abitur 2026 GK) \\ x → +∞: f'(x) → \leerfeld ; x → −∞: f'(x) → \leerfeld
\teil Die Ableitung einer Funktion $f$ ist $f'(x) = -x^2 \cdot (2x - 3) = -2x^3 + 3x^2$. Gib das Verhalten von $f'$ für $x \to -\infty$ und $x \to +\infty$ an. \hfill (Abitur 2026 GK) \\ x → +∞: f'(x) → \leerfeld ; x → −∞: f'(x) → \leerfeld
\teil $f(x) = -\frac{1}{4}x^4 + 2x$ hat die Ableitung $f'(x) = -x^3 + 2$. Gib das Verhalten von $f'$ für $x \to -\infty$ und $x \to +\infty$ an. \hfill (Abitur 2026 GK) \\ x → +∞: f'(x) → \leerfeld ; x → −∞: f'(x) → \leerfeld
\end{teile}
\end{aufgabe}

% Anlauf (ohne Prüfkennung)
\begin{aufgabe}{Produkte aus Polynom und e-Funktion – Anlauf}
\begin{teile}
\teil Welche Seite? Für welche Richtung geht bei $f(x) = (x - 6) \cdot e^{x}$ der e-Faktor gegen null? \kreuz{für $x \to +\infty$} \\ \kreuz{für $x \to -\infty$}
\teil $f(x) = e^{3x}$. Verhalten für $x \to +\infty$ und $x \to -\infty$? x → +∞: f(x) → \leerfeld ; x → −∞: f(x) → \leerfeld
\end{teile}
\end{aufgabe}

\begin{aufgabe}{Produkte aus Polynom und e-Funktion – Prüfungsaufgaben}
\begin{teile}
\teil $f(x) = (4 - x) \cdot e^{x}$. Gib das Verhalten für $x \to -\infty$ und $x \to +\infty$ an und sage, von welcher Seite sich der Graph der $x$-Achse nähert. \hfill (Abitur 2025 GK) \\ x → +∞: f(x) → \leerfeld ; x → −∞: f(x) → \leerfeld
\teil $f(x) = (x + 6) \cdot e^{x}$. Gib das Verhalten für $x \to -\infty$ und $x \to +\infty$ an und sage, von welcher Seite sich der Graph der $x$-Achse nähert. \hfill (Abitur 2025 GK) \\ x → +∞: f(x) → \leerfeld ; x → −∞: f(x) → \leerfeld
\teil $f(x) = (1 - 2x) \cdot e^{x}$. Gib das Verhalten für $x \to -\infty$ und $x \to +\infty$ an und sage, von welcher Seite sich der Graph der $x$-Achse nähert. \hfill (Abitur 2025 GK) \\ x → +∞: f(x) → \leerfeld ; x → −∞: f(x) → \leerfeld
\teil $f(x) = (x + 4) \cdot e^{-x}$. Grenzwert für $x \to +\infty$? Beschreibe den Verlauf des Graphen dort. \hfill (Abitur 2022 GK) \\ Grenzwert: \leerfeld
\teil $f(x) = (2x - 1) \cdot e^{-0{,}5x}$. Verhalten für $x \to +\infty$? Beschreibe den Verlauf des Graphen dort. \hfill (Abitur 2018 GK) \\ Grenzwert: \leerfeld
\teil $f(x) = (3x + 9) \cdot e^{-x}$. Grenzwert für $x \to +\infty$? Beschreibe den Verlauf des Graphen dort. \hfill (Abitur 2022 GK) \\ Grenzwert: \leerfeld
\teil $f(x) = 0{,}5 \cdot (x^2 - 9) \cdot e^{x}$. Verhalten für $x \to +\infty$ und $x \to -\infty$? \hfill (Abitur 2023 GK) \\ x → +∞: f(x) → \leerfeld ; x → −∞: f(x) → \leerfeld
\teil $f(x) = (-0{,}2x^2 + x) \cdot e^{-0{,}5x}$. Verhalten für $x \to +\infty$ und $x \to -\infty$? \hfill (Abitur 2021 GK) \\ x → +∞: f(x) → \leerfeld ; x → −∞: f(x) → \leerfeld
\teil $f(x) = 0{,}2x \cdot (4 - x) \cdot e^{x}$. Beschreibe den Verlauf des Graphen für $x \to -\infty$ und begründe am Term. \hfill (Abitur 2020 GK)
\teil Nach einem Regenguss ist der Pegel eines Bachs $w(t) = 6t \cdot e^{-0{,}5t} + 40$ ($t$ in h, $w$ in cm). Wie verhalten sich die Werte für $t \to +\infty$, und was bedeutet das? \hfill (Abitur 2019 GK)
\teil Nach einer Mahlzeit ist der Blutzucker $b(t) = 30t \cdot e^{-t} + 90$ ($t$ in h, $b$ in mg/dl). Wie verhalten sich die Werte für $t \to +\infty$, und was bedeutet das? \hfill (Abitur 2019 GK)
\teil Nach dem Einschalten einer Heizung ist die Temperatur eines Rohrs $T(t) = 12t \cdot e^{-0{,}2t} + 18$ ($t$ in min, $T$ in °C). Wie verhalten sich die Werte für $t \to +\infty$, und was bedeutet das? \hfill (Abitur 2019 GK)
\end{teile}
\end{aufgabe}

\begin{aufgabe}{Produkte aus Polynom und e-Funktion – Prüfungsaufgaben – weiter}
\begin{teile}
\teil $f(x) = (x - 5) \cdot e^{0{,}5x}$. Nullstelle und Verhalten für $x \to -\infty$ und $x \to +\infty$? \hfill (Abitur 2025 GK) \\ x = \leerfeld ; x → +∞: f(x) → \leerfeld ; x → −∞: f(x) → \leerfeld
\teil $f(x) = (6 - 2x) \cdot e^{x}$. Nullstelle und Verhalten für $x \to -\infty$ und $x \to +\infty$? \hfill (Abitur 2025 GK) \\ x = \leerfeld ; x → +∞: f(x) → \leerfeld ; x → −∞: f(x) → \leerfeld
\teil $f(x) = (x + 7) \cdot e^{-0{,}5x}$. Nullstelle und Verhalten für $x \to -\infty$ und $x \to +\infty$? \hfill (Abitur 2018 GK) \\ x = \leerfeld ; x → +∞: f(x) → \leerfeld ; x → −∞: f(x) → \leerfeld
\end{teile}
\end{aufgabe}
